\chapter{Experimental Results}
\label{chap:results}

This chapter validates the theoretical predictions of
Chapters~\ref{chap:irls}--\ref{chap:comparison} and puts both algorithms
side by side to determine when each one is preferable.
All experiments use the synthetic generator of Section~\ref{sec:impl_data};
scripts are in \texttt{experiments/}.

% ------------------------------------------------------------------
\section{Setup}
\label{sec:results_setup}

\begin{center}
\begin{tabular}{ll}
\toprule
\textbf{Parameter} & \textbf{Value} \\
\midrule
Problem size         & $n = 100$ features, $M = 300$--$400$ samples \\
True sparsity        & $10\%$ non-zero components \\
Noise / regulariser  & $\sigma = 0.05$,\quad $\lambda = 0.1$ \\
IRLS settings        & $\varepsilon_{\mathrm{thr}}=10^{-8}$, $\varepsilon_{\mathrm{stop}}=10^{-10}$,
                       $k_{\max}=100$, Cholesky \\
DSM settings         & $\beta=1.0$, $\delta_0=0.1f^*$, $\rho=0.95$, $i_{\max}=8\,000$ \\
Reference optimum    & scikit-learn \texttt{Lasso} at tolerance $10^{-12}$ \\
\bottomrule
\end{tabular}
\end{center}

% ------------------------------------------------------------------
\section{Convergence: a direct comparison}
\label{sec:results_convergence}

Chapters~\ref{chap:irls} and~\ref{chap:dsm} established the convergence
properties of each algorithm individually.
Here we confront them on the same problem.
Table~\ref{tab:conv_summary} recaps the key properties and the
empirically measured outcomes; Figures~\ref{fig:conv_iter}
and~\ref{fig:conv_time} show the gap curves.

\begin{table}[h!]
\centering
\begin{tabular}{lll}
\toprule
\textbf{Property} & \textbf{IRLS} & \textbf{DSM} \\
\midrule
Monotone iterates & \textbf{Yes} (MM guarantee) & No — record $\bar{f}^{\,i}$ only \\
Convergence rate  & \textbf{Linear} (geometric) & Sublinear $O(1/\varepsilon^2)$ \\
Gap after 100 iters   & $\mathbf{4.6 \times 10^{-7}}$ & $\approx 7.3 \times 10^{-2}$ \\
Gap after 8\,000 iters & N/A ($k_{\max}=100$) & $7.3 \times 10^{-2}$ (unchanged) \\
Accuracy advantage & \textbf{IRLS by $\sim 5$ orders of magnitude} & \\
\bottomrule
\end{tabular}
\caption{Convergence comparison on the base problem ($n=100$, $M=300$, $\lambda=0.1$).}
\label{tab:conv_summary}
\end{table}

\begin{figure}[h!]
    \centering
    \includegraphics[width=\textwidth]{figures/convergence_vs_iter.pdf}
    \caption{Optimality gap vs.\ iterations (semilog).
    IRLS (left) converges linearly and reaches machine-precision accuracy in under
    100 iterations. DSM (right) stagnates: the record $\bar{f}^{\,i} - f^*$ is
    orders of magnitude larger and barely moves after a few hundred steps.}
    \label{fig:conv_iter}
\end{figure}

\begin{figure}[h!]
    \centering
    \includegraphics[width=0.8\textwidth]{figures/convergence_vs_time.pdf}
    \caption{Gap vs.\ CPU time for both algorithms.
    IRLS reaches $10^{-6}$ accuracy in under 5\,ms; DSM never reaches $10^{-2}$
    even after several seconds, confirming that IRLS dominates on both axes
    in this regime ($n \le 100$).}
    \label{fig:conv_time}
\end{figure}

\paragraph{Non-monotonicity of DSM.}
Figure~\ref{fig:dsm_nonmonotone} highlights a practical consequence that
Table~\ref{tab:conv_summary} cannot fully convey: the raw iterates
$f(\mathbf{w}_i)$ oscillate wildly, while the record $\bar{f}^{\,i}$
decreases only occasionally.
This means that \emph{the last iterate is not the best one}; the algorithm
must track $\bar{f}^{\,i}$ and $\mathbf{w}_{\mathrm{best}}$ explicitly.
For IRLS no such bookkeeping is needed — every step is an improvement.

\begin{figure}[h!]
    \centering
    \includegraphics[width=\textwidth]{figures/dsm_nonmonotone.pdf}
    \caption{DSM raw iterates $f(w_i)$ (red) vs.\ monotone record $\bar{f}^{\,i}$ (black)
    vs.\ $f^*$ (green dashed). The gap between the two curves illustrates how
    misleading it would be to judge DSM's progress by the current iterate alone.}
    \label{fig:dsm_nonmonotone}
\end{figure}

\paragraph{Take-away.}
In the typical ELM training regime ($n \lesssim 500$, high accuracy required),
\textbf{IRLS is unambiguously superior}: it converges faster, monotonically, and
to far higher accuracy with fewer iterations, fewer parameters to tune, and a
reliable stopping criterion.
DSM's non-monotone, sublinear convergence is not a bug — it is the price of
handling non-smoothness without solving a linear system — but it makes DSM
unsuitable whenever high accuracy or exact sparsity is needed.

% ------------------------------------------------------------------
\section{Scalability: when does DSM catch up?}
\label{sec:results_scalability}

The convergence comparison above is decisive for $n=100$; the picture changes
as $n$ grows, because IRLS's Cholesky cost $O(n^3)$ eventually dominates.
We vary $n \in \{50, 100, 500, 1000, 3000\}$ with $M = 5n$
(Figure~\ref{fig:scalability}).

\begin{figure}[h!]
    \centering
    \includegraphics[width=\textwidth]{figures/scalability.pdf}
    \caption{Left: total wall-clock time vs.\ $n$ (log-log).
    Right: time per iteration vs.\ $n$.
    IRLS grows as $O(n^3)$ (Cholesky), DSM as $O(n^2)$ (since $M=5n$, per-iteration
    cost is $O(Mn)=O(n^2)$). The two curves cross near $n \approx 1000$.}
    \label{fig:scalability}
\end{figure}

The measured slopes confirm the theoretical cost model (Section~\ref{sec:comp_cost}):
IRLS scales as $O(n^3)$, DSM as $O(n^2)$.
They cross near $n \approx 1000$--$3000$, but the crossing is \emph{not enough
to declare DSM the winner}: at $n = 3000$, IRLS still achieves gap $\approx 2 \times 10^{-5}$
while DSM reaches only $\approx 1.6$ — three orders of magnitude worse in accuracy
despite the smaller wall-clock time.

\begin{center}
\begin{tabular}{ccccc}
\toprule
$n$ & IRLS time (s) & IRLS gap & DSM time (s) & DSM gap \\
\midrule
50   & 0.003 & $3.6 \times 10^{-7}$ & 0.06  & $3.4 \times 10^{-2}$ \\
100  & 0.005 & $6.0 \times 10^{-7}$ & 0.09  & $5.2 \times 10^{-2}$ \\
500  & 0.052 & $3.4 \times 10^{-6}$ & 1.16  & $3.6 \times 10^{-1}$ \\
1000 & 0.298 & $8.7 \times 10^{-6}$ & 6.78  & $5.4 \times 10^{-1}$ \\
3000 & 8.65  & $2.2 \times 10^{-5}$ & 61.3  & $1.6$ \\
\bottomrule
\end{tabular}
\end{center}

\paragraph{Interpretation.}
IRLS is faster \emph{and} more accurate at every scale tested.
DSM becomes \emph{time-competitive} only beyond $n \approx 1000$, but
its accuracy advantage never appears — its gap grows with $n$, while IRLS's gap
barely moves.
The only scenario where DSM would be genuinely preferable is one where
$n$ is so large that storing $\mathbf{X}^T\mathbf{X} \in \mathbb{R}^{n \times n}$
is impractical, or where $O(n^3)$ Cholesky is computationally prohibitive and
the CG inner solver of IRLS is insufficient — a regime beyond the scope of
this project's ELM setting.

% ------------------------------------------------------------------
\section{Head-to-head comparison}
\label{sec:results_comparison}

Table~\ref{tab:comparison} delivers the clearest verdict: for a problem with
$n = 100$, $M = 400$, we record the iteration count and CPU time for IRLS
to reach several accuracy thresholds, and whether DSM ever reaches the same
threshold within 10\,000 iterations.

\begin{table}[h!]
\centering
\begin{tabular}{ccccc}
\toprule
$\varepsilon = f - f^*$ & IRLS iters & IRLS time (s) & DSM iters & DSM time (s) \\
\midrule
$10^{-2}$ & \textbf{3}   & \textbf{0.0003} & $>10\,000$ & N/A \\
$10^{-3}$ & \textbf{8}   & \textbf{0.0005} & $>10\,000$ & N/A \\
$10^{-4}$ & \textbf{21}  & \textbf{0.0011} & $>10\,000$ & N/A \\
$10^{-6}$ & \textbf{113} & \textbf{0.0054} & $>10\,000$ & N/A \\
\bottomrule
\end{tabular}
\caption{Iterations and CPU time to first reach gap $\varepsilon$ ($n=100$, $M=400$,
$\lambda=0.1$). DSM never reaches any threshold within 10\,000 iterations; its
best recorded gap is $\approx 0.073$.}
\label{tab:comparison}
\end{table}

IRLS reaches $\varepsilon = 10^{-6}$ in 113 iterations and 5\,ms.
DSM, after 10\,000 iterations, has not reached even $\varepsilon = 10^{-2}$.
Solution quality confirms the gap: IRLS recovers the true sparsity pattern
(57\% zeros vs.\ 60\% true, $\|\mathbf{w}_{\mathrm{IRLS}} - \mathbf{w}^*\|
= 2.5 \times 10^{-4}$), while DSM produces a \emph{dense} solution
($\|\mathbf{w}_{\mathrm{DSM}} - \mathbf{w}^*\| = 0.23$, sparsity $= 0\%$).

\begin{figure}[h!]
    \centering
    \includegraphics[width=\textwidth]{figures/comparison_irls_dsm.pdf}
    \caption{IRLS vs.\ DSM on the same axes ($n=100$, $M=400$).
    Left: gap vs.\ iterations — IRLS's linear convergence is immediately visible
    as a straight line on the semilog plot, while DSM plateaus.
    Right: gap vs.\ CPU time — IRLS is faster by $3$--$4$ orders of magnitude
    at any given accuracy level.}
    \label{fig:comparison}
\end{figure}

\paragraph{When to prefer DSM.}
Despite losing on every quantitative metric in this setting, DSM retains
genuine practical value in scenarios where the Cholesky linear system solve is
infeasible: very large $n$ ($n \gtrsim 3000$) where IRLS's $O(n^3)$ cost
becomes prohibitive, memory-constrained environments where storing
$\mathbf{X}^T\mathbf{X}$ is impractical, or online/streaming settings where
$\mathbf{X}$ changes and precomputation cannot be amortised.
In all such cases, DSM provides a convergence guarantee (to $f^*$, albeit
sublinearly) with a per-iteration cost of only $O(Mn)$ and no matrix storage
beyond the current iterate.
